Knowledge graphs have long ago become an attractive representation for structured knowledge extracted
from text. Early work on ontology learning established methods for inducing concept hierarchies
and relations from corpora: \citet{hearst1992hyponyms} introduced lexico-syntactic
patterns for hyponymy acquisition, while \citet{maedche2001ontologylearning} combined
linguistic and statistical co-occurrence techniques in a semiautomatic ontology
engineering framework. \citet{hogan2021kgsurvey} provide a comprehensive
overview of KG construction, population, and enrichment approaches spanning rule-based
systems, distant supervision, and embedding-based methods.

The advent of large language models has substantially lowered the barrier to KG
construction at scale. Tools such as LlamaIndex~\cite{liu2022llamaindex},
LangChain~\cite{chase2022langchain}, and the Neo4j LLM Knowledge Graph
Builder~\cite{neo4j2024llmgraphbuilder} expose pipelines in which an LLM is prompted to
extract entities and relations from text, optionally constrained by a user-supplied schema.
These tools have seen wide adoption but share a fundamental design choice: the ontology
--- node types, relation types, and constraints --- is specified once and remains fixed for
the duration of construction.

More recent academic approaches address extraction quality without revisiting the ontology
question. \citet{mo2025kggen} propose KGGen, which extracts KGs from plain text using
language models with improved entity disambiguation. \citet{tsang2025autographr1} apply
reinforcement learning to optimise KG construction end-to-end, using graph quality signals
as reward. \citet{mynarz2023testdriven} propose a test-driven approach in which competency
questions drive extraction, providing a form of task specification without a quantitative
feedback loop. In all these methods, ontology design remains a pre-construction decision
rather than an adaptive variable.

\citet{li2025llmontologyreview} provide the most comprehensive survey to date of LLMs
across the ontology engineering lifecycle, reviewing 30 publications and 41 experiments
spanning requirements specification, conceptualisation, encoding, ontology matching, and
maintenance. Their analysis reveals a pronounced imbalance: ontology implementation tasks
account for 63.4\% of studied experiments, while maintenance has been addressed by a single
study (2.4\%). Crucially, no reviewed approach evaluates ontology quality through downstream
task performance; evaluation in all cases targets structural properties --- axiom soundness,
semantic consistency, or alignment with competency questions. The survey identifies as an
open direction ``developing hybrid learning pipelines that allow LLMs to continuously
integrate domain updates\ldots\ thereby supporting ongoing validation and refinement based
on domain-specific feedback''~\citep{li2025llmontologyreview}.

The evaluation of knowledge graphs has historically been conducted through intrinsic,
structural quality metrics. \citet{paulheim2017kgquality} survey a broad class of such
metrics, covering completeness, consistency, and conciseness. More recently,
\citet{paulheim2024structuralquality} draw an explicit distinction between this intrinsic
mode of evaluation and task-based evaluation --- assessing a KG by its utility on
a concrete downstream application. This distinction is central to the motivation of the
present work: a graph may score well on structural criteria while failing to support a
specific predictive task.

Several frameworks have operationalised task-based KG evaluation. KGBench~\cite{bloem2021kgbench}
provides a suite of relational and multimodal machine learning tasks benchmarked on
curated KG datasets. GEval~\cite{choi2020geval} offers a software framework for evaluating
KGs through downstream task performance across multiple dimensions.
KGrEaT~\cite{heist2023kgreat} enables systematic comparison of KG variants across
classification, regression, and recommendation tasks; the framework explicitly supports
comparison of ``different versions of a KG during its life cycle,'' precisely the kind of
comparison our feedback loop performs automatically. LLM-KG-Bench~\cite{aksw2024llmkgbench}
focuses specifically on the quality of LLM-generated graphs. Together, these frameworks
confirm that task utility cannot be inferred from structural properties alone --- yet none
of them feeds performance signals back into the construction process itself.

The closest prior work to our approach concerns task-aware graph construction.
\citet{cucumides2025augraph} propose constructing graphs from tabular and relational data
to maximise GNN performance on a target task, treating graph topology as a
task-dependent variable optimised through downstream loss. However, their setting
involves structured tabular input rather than unstructured text, and the construction is
not driven by an evolving ontology.

In the financial domain, \citet{arun2025finreflectkg} introduce FinReflectKG, an agentic
framework for constructing and evaluating financial knowledge graphs. Their system includes
an evaluation component that scores the extracted KG, but feedback is qualitative and does
not close a quantitative loop between classifier performance and ontology revision. Both
works share the recognition that graph structure should be task-conditioned; neither
implements an automated, iterative feedback mechanism of the kind we propose.

A complementary body of work addresses how to derive informative features from knowledge
graphs for downstream prediction. Metapath-based approaches~\cite{sun2011pathsim,
dong2017metapath2vec} exploit typed paths in heterogeneous graphs to capture relational
semantics between node pairs. Relational Graph Convolutional Networks
(R-GCN)~\cite{schlichtkrull2018rgcn} extend GCNs to multi-relational graphs, learning
embeddings that respect relation type. Knowledge graph embedding methods survey a broad
class of techniques for representing entities and relations in continuous
space~\cite{wang2017kgsurvey, ji2021kgcompletion}, supporting link prediction and
node-level classification.

These methods share the assumption that the knowledge graph is given and fixed; feature
engineering, embedding, or message-passing then operates over the resulting structure. In
our framework we draw on metapath counting and relation type histograms as components of
the feature vector, but the graph from which these features are derived is itself adapted
in response to the downstream task. Rather than learning better from a fixed
graph, we ask what graph structure produces the most learnable features.

Taken together, the literature reveals a consistent gap. Evaluation frameworks demonstrate
that KG utility is task-dependent~\cite{heist2023kgreat, bloem2021kgbench}; task-aware
construction methods show that graph structure can be conditioned on downstream
objectives~\cite{cucumides2025augraph}; and graph feature methods show how to extract
predictive signal from heterogeneous graphs~\cite{sun2011pathsim, schlichtkrull2018rgcn}.
Yet no existing work combines these strands into a closed feedback loop in which a
lightweight node-level predictor directly steers ontology evolution at construction time.
The present paper addresses this gap by treating the KG ontology as a hyperparameter
and iteratively refining it based on quantitative downstream performance signals.
